\chapter{Состав набора для сборки компьютера}

В таблице указано максимальное число модулей, которого будет достаточно
для выполнения любого из заданий.

\begin{tabular}{|l|c|}
\hline
\bf Модуль & \bf Количество \\ \hline
Переключатели & \theswitches \\ \hline
Регистр & \theregister \\ \hline
Унарные операции & \theunary \\ \hline
Бинарные операции & \thebinary \\ \hline
Сумматор & \theadder \\ \hline
Дешифратор & \thedecoder \\ \hline
Шина & \thebus \\ \hline
Кросс & \thematrix \\ \hline
Тактовый генератор & \theclock \\ \hline
ПЗУ & \therom \\ \hline
Адаптер блока питания & \thepower \\ \hline
\end{tabular}

Для сборки всех этих модулей потребуется соответствующее число плат и следующие компоненты.

\the\numexpr \theadder * 8 + \theregister * 8 \relax

\begin{tabular}{|l|c|}
\hline
\bf Компонент & \bf Количество \\ \hline
\input{appendix/bom.tex}
\unskip блок питания & 1 \\ \hline
соединительный шлейф & \thecable \\ \hline
\end{tabular}

Угловые и вертикальные пины могут быть собраны в линейки разной длины.
Для сборки их необходимо разрезать самостоятельно.

Кроме перечисленных компонентов, также понадобится некоторое количество
перемычек и проводов для соединения входов и выходов в кросс-модулях.
